%% =============================================================================
%% PAPER 1 - Section 4: Benchmark Factor Models
%% =============================================================================
%% Full prose. Tables/figures referenced via \ref{} only.
%% Actual tables in paper1_tables_sec4.tex; figures in paper1_figures_sec4.tex.
%% Maps to: 02a, 02b, 02c, 03_subperiod_rolling_analysis
%% =============================================================================


\subsection{Benchmark Specifications}
\label{sec:bench_specs}

% TODO: Full prose covering:
%
% Three canonical factor models:
%   Duarte et al. (2007): 10 factors (Mkt-RF, SMB, HML, UMD, RS, RI, RB, R2, R5, R10)
%   Fung & Hsieh (2004): 7 factors (SNP, SIZE, PTFSBD, PTFSFX, PTFSCOM, TERM, CREDIT)
%   Brooks et al. (2020): 8 Active FI premia (Term, Global_Term, ..., UST_Volatility)
%
% All estimated via OLS with Newey-West HAC standard errors (6 lags).
% EUR factors used for primary specification.
% Monthly frequency.


\subsection{Results and Discussion}
\label{sec:bench_results}

% TODO: Full prose covering:
%
% Factor model regressions:
%   Alpha significant across all 3 frameworks and all 3 strategies
%   R-squared low (5-15%): benchmarks do not span returns
%   Factor loadings discussed
%   → Tables \ref{tab:duarte_article}, \ref{tab:funghsieh_article}, \ref{tab:activefi_article}
%
% Cross-model comparison and regime analysis:
%   Alpha robust across models (Table \ref{tab:alpha_synthesis})
%   Alpha higher during stress (Panel B: HIGH vs NORMAL)
%   Consistent with Duffie (2010): limits to arbitrage, slow-moving capital
%   Motivates cross-strategy analysis in Section 6
%
% Rolling alpha:
%   Visual confirmation: alpha spikes during stress episodes
%   Pattern holds across all 3 benchmark specifications
%   → Figure \ref{fig:rolling_alpha_composite}
%
% Appendix \ref{app:subperiod}: sub-period splits, threshold robustness,
%   per-framework rolling alpha plots